% !TEX root =  main.tex
%%%%%%%%%%%%%%%%%%%%%%%%%%%%%%%%%%%%%
%%%%%%%%%%%%%%%%%%%%%%%%%%%%%%%%%%%%%

\begin{abstract}
    % As large language models (LLMs) become increasingly deployed in critical applications, ensuring their safety and robustness through effective red-teaming has become paramount. While recent automated red-teaming methods show promise, most existing approaches rely on fixed, human-specified red-teaming frameworks, optimizing attack policies or strategies within predefined interaction loops. This reliance on manually designed workflows suffer from human biases and make exploring the broader design expensive. We introduce \method, an automated pipeline that leverages LLMs’ in-context learning capability to iteratively design and refine red-teaming systems without human intervention. Rather than optimizing attacker policies within a predefined interaction structure, \method~is the first work to treat red-teaming as a system design problem and apply evolutionary algorithms to automatically optimize the workflow of red-teaming systems. Inspired by Meta Agent Search, a paradigm where LLMs iteratively design and refine other agentic systems, we tailor \method~specifically for the red-teaming domain, improving system performance through evolutionary selection. We demonstrate that red-teaming systems designed by \method~consistently outperform state-of-the-art approaches, achieving 96\% attack success rate (ASR) on Llama-2-7B (36\% improvement) and 98\% (10\% improvement) on Llama-3-8B on HarmBench. Our approach exhibits strong transferability to proprietary models, achieving 100\% ASR on GPT-3.5-Turbo and GPT-4o, and 60\% on Claude-Sonnet-3.5 (24\% improvement). This work highlights automated agentic system design as a powerful paradigm for robust AI safety evaluation that can keep pace with rapidly evolving models.
While recent automated red-teaming methods show promise for systematically exposing model vulnerabilities, most existing approaches rely on human-specified workflows. This dependence on manually designed workflows suffers from human biases and makes exploring the broader design space expensive. We introduce \method, an automated pipeline that leverages LLMs' in-context learning to iteratively design and refine red-teaming systems without human intervention. Rather than optimizing attacker policies within predefined structures, \method~treats red-teaming as a system design problem, and it autonomously evolves automated red-teaming systems using evolutionary selection and generational knowledge. Red-teaming systems designed by \method~consistently outperform state-of-the-art approaches, achieving 96\% attack success rate (ASR) on Llama-2-7B, 98\% on Llama-3-8B and 100\% on Qwen3-8B on HarmBench. Our approach generates robust, query-agnostic red-teaming systems that transfer strongly to the latest proprietary models, achieving an impressive 100\% ASR on GPT-5.1, DeepSeek-R1 and DeepSeek V3.2. This work highlights evolutionary algorithms as a powerful approach to AI safety that can keep pace with rapidly evolving models.
\end{abstract}